% ================================================================
% APPENDIX: Assignment and Rubric Reference
% EDNS 491/492 Capstone Design I & II
% ================================================================

\newpage
\appendix
\section{Assignment and Rubric Reference}
\label{sec:assignments-rubrics}

This appendix consolidates the key assignments and rubrics for both EDNS~491 (Senior Design~I) and EDNS~492 (Senior Design~II). Use this as a quick reference for understanding expectations, point allocations, and grading criteria. \textbf{Canvas remains the authoritative source} for due dates, submission links, and any mid-semester updates to rubrics.

\begin{warnbox}
Rubric criteria summaries in this appendix are condensed for reference. Always review the full rubric on Canvas before submitting---some criteria contain additional detail not captured here.
\end{warnbox}

% ──────────────────────────────────────────────────────────────────
\subsection{EDNS 491 --- Senior Design I Assignments}
\label{subsec:sdi-assignments}
% ──────────────────────────────────────────────────────────────────

\subsubsection{Sprint Cycle Assignments (Sprints 0--5)}
\label{sssec:sprint-cycles-sdi}

Sprint Cycles form the backbone of your Scrum-based project management. Each two-to-three-week sprint follows the same structure:

\begin{enumerate}[leftmargin=*]
  \item \textbf{Sprint Planning} --- Define sprint goals, populate the Sprint Backlog, assign tasks.
  \item \textbf{Execution} --- Daily stand-ups, design work, client engagement.
  \item \textbf{Sprint Review} --- Present progress to client and PA (slide deck required).
  \item \textbf{Sprint Retrospective} --- Team reflection on process improvements (\emph{kaizen}).
\end{enumerate}

\noindent\textbf{Submission (each sprint):} Upload (1) the Sprint Review slide deck and (2) the team retrospective within 24 hours of your Sprint Review meeting.

\begin{tipbox}
Sprint Cycles 0--5 are scored at 0 points individually but contribute to your \emph{Team Assessment -- Project Progress} grade (10\% of course). Your PA evaluates sprint engagement holistically.
\end{tipbox}

\noindent The sprint themes in SDI are:
\begin{itemize}[leftmargin=*]
  \item \textbf{Sprint 0} --- Engaging with the Problem (client kickoff preparation and meeting)
  \item \textbf{Sprint 1} --- Planning Through Ambiguity (90\% complete Project Plan)
  \item \textbf{Sprint 2} --- Trade Study \& Concept Development (concept selection table)
  \item \textbf{Sprint 3} --- Defining Your Concept (primary concept with client alignment)
  \item \textbf{Sprint 4} --- From Concept to Preliminary Design (prototyping, feasibility evidence)
  \item \textbf{Sprint 5} --- Preliminary Design Reviews (PDR presentation and report)
\end{itemize}

% ─── Project Kickoff ──────────────────────────────────────────────
\subsubsection{Project Kickoff \& Client Meeting}
\label{sssec:kickoff}

\begin{tabularx}{\textwidth}{lX}
\toprule
\textbf{Points} & 10 (team) \\
\textbf{Deliverables} & Team Information Slide, Stakeholder Document, Sprint~0 slide deck, follow-up email to client \\
\textbf{Key expectations} & Prepare a 5-slide deck demonstrating early engagement; send follow-up email within 24 hrs; document next steps and first Sprint Review date \\
\bottomrule
\end{tabularx}

\begin{protip}
Arrive at the client meeting with thoughtful, open-ended questions about the \emph{problem}, not premature solutions. Print color copies of your Team Information Slide for the client.
\end{protip}

% ─── Statement of Work / Project Plan ─────────────────────────────
\subsubsection{Project Plan (Statement of Work)}
\label{sssec:project-plan}

\begin{tabularx}{\textwidth}{lX}
\toprule
\textbf{Points} & 100 (team) --- initial draft is ungraded feedback; updated submission is graded \\
\textbf{Audience} & Client (primary), PA, Technical Advisors \\
\textbf{Key sections} & Project summary, scope of work, customer needs, design requirements (tabulated), stakeholders, background research, engineering standards, budget, material handoff, design process, project schedule (Gantt chart), team members, management tools, communication strategies \\
\bottomrule
\end{tabularx}

\begin{warnbox}
Do \textbf{not} send the initial draft submission to the client. The initial submission is for PA feedback only. The updated submission is shared with the client via email.
\end{warnbox}

% ─── PDR Report ───────────────────────────────────────────────────
\subsubsection{Preliminary Design Report (PD Report)}
\label{sssec:pdr-report}

\begin{tabularx}{\textwidth}{lX}
\toprule
\textbf{Points} & 100 (team) --- initial submission for PA feedback; updated submission is graded \\
\textbf{Audience} & Client (stand-alone document) \\
\textbf{Key sections} & Executive summary, scope \& background, design requirements, concept exploration (3+ concepts with diagrams), concept critique (sustainability/equity/risk tool), concept selection, functional/spatial basis, life cycle analysis, detailed drawings, budget, schedule \\
\bottomrule
\end{tabularx}

\begin{protip}
Submit the initial PD Report to your client \emph{before} the Preliminary Design Review so they can prepare questions. The updated version incorporates feedback from the review and is attached to the Next Steps Email.
\end{protip}

% ─── PDR Presentation ─────────────────────────────────────────────
\subsubsection{Preliminary Design Review (Presentation)}
\label{sssec:pdr-presentation}

\begin{tabularx}{\textwidth}{lX}
\toprule
\textbf{Points} & 100 (team + individual components) \\
\textbf{Format} & 15--20 min presentation + 15--20 min Q\&A \\
\textbf{Audience} & Client, mentors, PA, faculty, students \\
\textbf{Requirements} & Use Capstone presentation template; every member speaks and is listed on their slides; business casual or formal attire required \\
\bottomrule
\end{tabularx}

% ─── Ethics ───────────────────────────────────────────────────────
\subsubsection{Ethics Assignment}
\label{sssec:ethics}

\begin{tabularx}{\textwidth}{lX}
\toprule
\textbf{Points} & 100 (individual) \\
\textbf{Format} & 750--1500 word essay, PDF, Times New Roman or Arial 12pt \\
\textbf{Task} & Identify a specific ethical dilemma in your Capstone project; analyze using engineering ethical frameworks and relevant Code of Ethics; determine a professional course of action \\
\textbf{Key distinction} & Must present a genuine \emph{dilemma} (competing values), not a generic safety statement \\
\bottomrule
\end{tabularx}

\noindent\textbf{Required sections:} (1)~The dilemma, (2)~Stakeholder analysis, (3)~Alternative courses of action, (4)~Application of Canons, (5)~Rationale \& decision, (6)~Measurement, (7)~References (IEEE), (8)~Optional Appendix~A if AI was used.

\begin{warnbox}
AI may be used as a ``simulator and debater'' (stakeholder simulation, blind-spot checks, proofreading) but \textbf{not} to write the draft or generate citations. If AI is used, Appendix~A (AI Transparency Statement) is mandatory.
\end{warnbox}

% ─── Engineering Standards ────────────────────────────────────────
\subsubsection{Engineering Standards}
\label{sssec:eng-standards}

\begin{tabularx}{\textwidth}{lX}
\toprule
\textbf{Points} & 100 (individual) \\
\textbf{Task} & Select a relevant standard/regulation module, summarize it, explain how it applies to your project, share findings with the team, and schedule a team discussion with your PA \\
\textbf{Deliverables} & Summary document uploaded to Canvas and to the team's shared folder \\
\bottomrule
\end{tabularx}

% ─── Goal Setting / Attainment ────────────────────────────────────
\subsubsection{Project Goal Setting \& Goal Attainment}
\label{sssec:goals}

\textbf{Goal Setting} (10~pts, team): Early-semester document stating end-of-semester goals agreed upon by client and PA, with a clear definition of done.

\textbf{Goal Attainment} (100~pts, team): End-of-semester document restating goals and providing evidence of achievement (images, videos, links to drawings, etc.). PA evaluates based on reasonable goal-setting and systematic progress through sprint reviews.

\begin{warnbox}
If your team does not have a functional prototype, preliminary site drawings, schematics, or equivalent by the end of SDI, the team will not receive an A regardless of other assignment scores.
\end{warnbox}

% ─── Individual Performance Evaluation ────────────────────────────
\subsubsection{Individual Performance Evaluations}
\label{sssec:ipe-sdi}

\begin{tabularx}{\textwidth}{lX}
\toprule
\textbf{Eval \#1} & 100 pts, mid-semester \\
\textbf{Eval \#2} & 110 pts, end-of-semester (includes semester reflection questions) \\
\textbf{Assessed by} & Project Advisor, informed by peer evaluations, sprint review observations, and contribution to deliverables \\
\textbf{Submission} & Upload a short document describing your technical and professional contributions with evidence (links to documents, deliverables) \\
\bottomrule
\end{tabularx}

\begin{tipbox}
If your evaluation falls below 70\%, your PA will issue a Performance Improvement Plan (PIP). Address it promptly---this is a professional development tool, not a punishment.
\end{tipbox}

% ─── Next Steps Email ─────────────────────────────────────────────
\subsubsection{Next Steps Email}
\label{sssec:next-steps-email}

\begin{tabularx}{\textwidth}{lX}
\toprule
\textbf{Points} & 10 (team) \\
\textbf{Audience} & Client (CC your PA and Technical Advisors) \\
\textbf{Key elements} & Thank client, summarize PDR feedback and how it was addressed (action items table), engineering recommendation for next semester, proposed budget, SDII schedule, attach updated PD Report \\
\textbf{Submission} & Paste email text or state the date sent \\
\bottomrule
\end{tabularx}

% ──────────────────────────────────────────────────────────────────
\subsection{EDNS 492 --- Senior Design II Assignments}
\label{subsec:sdii-assignments}
% ──────────────────────────────────────────────────────────────────

\subsubsection{Sprint Cycle Assignments (Sprints 7--14)}
\label{sssec:sprint-cycles-sdii}

SDII continues the Scrum cadence with Sprints 7--14. The structure mirrors SDI (planning, execution, review, retrospective). Key sprint milestones in SDII:

\begin{itemize}[leftmargin=*]
  \item \textbf{Sprint 7} --- Re-engagement with client and team; action plan for semester
  \item \textbf{Sprint 8} --- Critical Design Review preparation and delivery
  \item \textbf{Sprint 9} --- Post-CDR iteration; first peer evaluation
  \item \textbf{Sprints 10--11} --- Build, test, redesign cycle; mid-semester evaluation
  \item \textbf{Sprint 12} --- Showcase preparation; ``What We Wish We Would Have Known'' video
  \item \textbf{Sprint 13} --- Final Design Review and Report
  \item \textbf{Sprint 14} --- Design Showcase and final deliverables
\end{itemize}

% ─── Re-engagement ────────────────────────────────────────────────
\subsubsection{Re-engagement Email \& Activity}
\label{sssec:reengagement}

\begin{tabularx}{\textwidth}{lX}
\toprule
\textbf{Email} & Professional re-engagement email to client at start of SDII \\
\textbf{Activity} & Group reflection on SDI and action plan for SDII improvement \\
\bottomrule
\end{tabularx}

% ─── CDR Report ───────────────────────────────────────────────────
\subsubsection{Critical Design Report (CDR)}
\label{sssec:cdr-report}

\begin{tabularx}{\textwidth}{lX}
\toprule
\textbf{Points} & 100 (team) \\
\textbf{Key sections} & Design narrative, engineering analysis \& calculation package (with TA sign-off), risk assessment (FMEA or risk matrix), test plans \& procedures (V\&V), technical drawings, technical writing \\
\textbf{Critical requirement} & All calculations must have sign-off from a Technical Advisor or faculty member \\
\bottomrule
\end{tabularx}

\begin{warnbox}
The CDR is your ``license to build.'' If your engineering analysis has gaps or lacks TA sign-off, you risk delays in moving to fabrication and testing.
\end{warnbox}

% ─── CDR Presentation ─────────────────────────────────────────────
\subsubsection{Critical Design Review (Presentation)}
\label{sssec:cdr-presentation}

\begin{tabularx}{\textwidth}{lX}
\toprule
\textbf{Points} & 100 (team + individual) \\
\textbf{Criteria} & Context (20\%), Content (25\%), Process (25\%), Presentation (20\%), Graphics (10\%) \\
\textbf{Format} & Same structure as PDR; use Capstone presentation template; business attire required \\
\bottomrule
\end{tabularx}

% ─── FDR Report ───────────────────────────────────────────────────
\subsubsection{Final Design Report (FDR)}
\label{sssec:fdr-report}

\begin{tabularx}{\textwidth}{lX}
\toprule
\textbf{Points} & 100 (team) \\
\textbf{Key sections} & Project overview, design approach, engineering analysis, final solution (referencing how design criteria were met or not), project management (budget, WBS, schedule), lessons learned, technical writing \\
\textbf{Purpose} & Stand-alone document enabling future engineers to continue the work \\
\bottomrule
\end{tabularx}

% ─── FDR Presentation ─────────────────────────────────────────────
\subsubsection{Final Design Review (Presentation)}
\label{sssec:fdr-presentation}

\begin{tabularx}{\textwidth}{lX}
\toprule
\textbf{Points} & 100 (team + individual) \\
\textbf{Criteria} & Context (20\%), Content (25\%), Process (25\%), Presentation (20\%), Graphics (10\%) \\
\textbf{Note} & Must be completed \emph{before} the Design Showcase \\
\bottomrule
\end{tabularx}

% ─── Broader Impacts Essay ────────────────────────────────────────
\subsubsection{Broader Impacts Essay}
\label{sssec:broader-impacts}

\begin{tabularx}{\textwidth}{lX}
\toprule
\textbf{Points} & 100 (individual) \\
\textbf{Format} & Argumentative essay, 1000--1500 words, use provided template \\
\textbf{Task} & Argue how an engineered system related to your project has positively or negatively impacted society, the environment, and/or the economy. Must take a clear thesis and defend it against opposing arguments. \\
\textbf{Key distinction} & This is an \emph{argument}, not a report. Must include opposing viewpoints and rebuttals. \\
\bottomrule
\end{tabularx}

\noindent\textbf{Required sections:} (1)~Introduction \& thesis, (2)~Background, (3)~Proponent arguments, (4)~Opponent arguments, (5)~Rebuttal, (6)~Conclusion, (7)~References (IEEE), (8)~Optional Appendix~A if AI was used.

\begin{tipbox}
The top 10 essays are sent to a panel of alumni and faculty judges. Winners are recognized at the Awards Breakfast and published in the Design Showcase program.
\end{tipbox}

% ─── Design Showcase ──────────────────────────────────────────────
\subsubsection{Design Showcase (Poster)}
\label{sssec:showcase}

\begin{tabularx}{\textwidth}{lX}
\toprule
\textbf{Poster Draft} & 25 pts (team) --- submit well before printing for PA feedback \\
\textbf{Showcase} & 100 pts (team) --- scored by external judges (alumni, industry, faculty) \\
\textbf{Poster specs} & 48'' W $\times$ 36'' H in PowerPoint; min 24pt font in text boxes; 28--32pt recommended \\
\textbf{Judging criteria} & Poster \& display, project overview, design approach, design solution, next steps (each on 1--5 scale) \\
\textbf{Attendance} & Mandatory; business casual or professional recommended \\
\bottomrule
\end{tabularx}

% ─── Project Synopsis ─────────────────────────────────────────────
\subsubsection{Project Synopsis}
\label{sssec:synopsis}

\begin{tabularx}{\textwidth}{lX}
\toprule
\textbf{Points} & 10 (team) --- complete/incomplete grading \\
\textbf{Purpose} & Short summary for the Design Showcase program booklet \\
\textbf{Requirement} & Professional, concise, complete, fully proofread \\
\bottomrule
\end{tabularx}

% ─── Digital Design Archive ───────────────────────────────────────
\subsubsection{Digital Design Archive}
\label{sssec:archive}

End-of-project submission of all project files, documentation, CAD files, code, and deliverables in an organized digital archive for the client and future teams.

% ─── What We Wish We Would Have Known ─────────────────────────────
\subsubsection{``What We Wish We Would Have Known'' Video}
\label{sssec:wish-video}

Team-produced video sharing advice and lessons learned for future Capstone teams. Shown to incoming SDI students.

% ─── Manufacturing Readiness Review ───────────────────────────────
\subsubsection{Manufacturing Readiness Review}
\label{sssec:mrr}

Teams working in the Labriola truck bays or Thorson Design Center must complete safety training, hazard assessments, and access contracts before beginning fabrication. EHS sign-off is required for any safety-related work, chemical usage, or painting.


% ══════════════════════════════════════════════════════════════════
\subsection{Rubric Reference}
\label{subsec:rubric-reference}
% ══════════════════════════════════════════════════════════════════

The following tables summarize the grading criteria for the most important assignments. CLO (Course Learning Outcome) rows that do not count toward scores are omitted. Point values shown are the maximum for each criterion.

% ─── PDR Report Rubric ────────────────────────────────────────────
\subsubsection{Preliminary Design Report Rubric}
\label{sssec:rubric-pdr-report}

{\small
\begin{longtable}{p{3.2cm} c p{3.5cm} p{3.5cm} p{2.8cm}}
\toprule
\textbf{Criterion} & \textbf{Pts} & \textbf{Exceeds} & \textbf{Meets / Adequate} & \textbf{Marginal} \\
\midrule
\endfirsthead
\toprule
\textbf{Criterion} & \textbf{Pts} & \textbf{Exceeds} & \textbf{Meets / Adequate} & \textbf{Marginal} \\
\midrule
\endhead
\bottomrule
\endfoot

Executive Summary & 3 & Succinct, highlights all key points & May omit components or lack conciseness & Limited info; must read full report \\
\addlinespace
Project Scope \& Background & 5 & Thorough problem statement, existing solutions, standards & Generally clear but missing some depth & Insufficient background \\
\addlinespace
Design Requirements & 5 & Tabulated with verification methods, types, numbering & Table present but incomplete & Minimal discussion \\
\addlinespace
Concept Exploration & 40 & 3+ unique, feasible concepts with diagrams, functional/spatial basis, pros/cons & Fewer concepts or missing functional basis & One idea with little exploration \\
\addlinespace
Concept Critique & 15 & Actionable design updates from appropriate sustainability/equity/risk tool & Tool applied but results not well summarized & Poor attempt or no tool applied \\
\addlinespace
Concept Selection & 10 & Analysis-driven recommendation linked to constraints & Some reasoning but gaps in justification & Unsupported opinion \\
\addlinespace
Detailed Drawings & 4 & Complete, professional set in appendices & Present but incomplete or unprofessional & Not complete nor professional \\
\addlinespace
Budget & 3 & Discusses allocation, BoM, initial expenses & Present but lacking detail & Missing \\
\addlinespace
Project Schedule & 3 & Updated Gantt chart referenced in text, milestones discussed & Present but not fully discussed & Only included, not discussed \\
\addlinespace
Technical Writing & 10 & Clear, concise, one voice, proofread & Generally well written with some errors & Many errors; difficult to read \\
\addlinespace
Formatting \& Citations & 2 & IEEE citations, ToC, consistent formatting & Some formatting or citation issues & No citations; poorly formatted \\

\end{longtable}
}

% ─── PDR Presentation Rubric ─────────────────────────────────────
\subsubsection{Preliminary Design Review Presentation Rubric}
\label{sssec:rubric-pdr-pres}

{\small
\begin{longtable}{p{3.5cm} c p{4cm} p{4cm}}
\toprule
\textbf{Criterion} & \textbf{Pts} & \textbf{Exceeds Mastery} & \textbf{Below Mastery} \\
\midrule
\endfirsthead
\toprule
\textbf{Criterion} & \textbf{Pts} & \textbf{Exceeds Mastery} & \textbf{Below Mastery} \\
\midrule
\endhead
\bottomrule
\endfoot

Design Considerations & 16 & Thorough problem definition; precise goals, constraints, requirements; strong justification of safety, standards, broader impacts & Vague problem/goals; unclear requirements; little attention to safety or standards \\
\addlinespace
Proposed Solutions & 20 & Advanced engineering concepts with strong analysis and clear documentation; deep understanding of tradeoffs & Concepts applied incorrectly; minimal analysis; little awareness of tradeoffs \\
\addlinespace
Engineering Process & 20 & Sophisticated concepts, in-depth analysis, comprehensive documentation & Underdeveloped concepts; weak analysis \\
\addlinespace
Team Presentation & 30 & Exceptionally clear, well-organized, engaging; polished delivery & Lacks clarity or organization; unprepared delivery \\
\addlinespace
Communication \& Graphics & 10 & Original, compelling graphics; proper citations; minimal text & Ineffective graphics; excessive text; no citations \\
\addlinespace
Professional Attire & 4 & All members in business casual/formal & Lack of professional attire detracts from credibility \\

\end{longtable}
}

% ─── CDR Report Rubric ────────────────────────────────────────────
\subsubsection{Critical Design Report Rubric}
\label{sssec:rubric-cdr-report}

{\small
\begin{longtable}{p{3.2cm} c p{3.5cm} p{3.5cm} p{2.8cm}}
\toprule
\textbf{Criterion} & \textbf{Pts} & \textbf{Exceeds} & \textbf{Meets / Adequate} & \textbf{Marginal} \\
\midrule
\endfirsthead
\toprule
\textbf{Criterion} & \textbf{Pts} & \textbf{Exceeds} & \textbf{Meets / Adequate} & \textbf{Marginal} \\
\midrule
\endhead
\bottomrule
\endfoot

Design Narrative & 10 & Clear method/logic; narrative guides reader; work concisely summarized & Present but lacking in some areas & Limited attempt; significant improvement needed \\
\addlinespace
Engineering Analysis \& Calcs & 20 & Exceptionally detailed, error-free, TA-approved calc package & Some gaps or questions unanswered & Incomplete or lacking TA sign-off \\
\addlinespace
Risk Assessment & 15 & Detailed critique with risk tool (FMEA/matrix); clear mitigation plans linked to top risks & Tool present but surface-level analysis & Feels like afterthought; no clear narrative \\
\addlinespace
Test Plans (V\&V) & 10 & Complete, descriptive; non-team-member could execute tests & Adequate detail but needs improvement & Work needs significant improvement \\
\addlinespace
Technical Drawings & 15 & Professional-quality, properly annotated, complete set & Present but formatting issues or incomplete & Limited attempt at professional quality \\
\addlinespace
Technical Knowledge & 5 & Appropriate for completion level; clear justification of calculations and design decisions & Adequate but needs improvement & Significant improvement needed \\
\addlinespace
Technical Writing & 15 & Clear, concise, one voice, proofread & Several errors affecting clarity & Many errors; not senior-level \\

\end{longtable}
}

% ─── CDR Presentation Rubric ─────────────────────────────────────
\subsubsection{Critical Design Review Presentation Rubric}
\label{sssec:rubric-cdr-pres}

{\small
\begin{longtable}{p{4cm} c p{5cm} p{3cm}}
\toprule
\textbf{Criterion} & \textbf{Pts} & \textbf{Full Marks} & \textbf{Poor} \\
\midrule
\endfirsthead
\bottomrule
\endfoot

Context & 20 & Succinctly defines problem; clear goals, constraints; addresses safety and broader impacts & 10 \\
\addlinespace
Content & 25 & Appropriate engineering concepts, analysis, documentation; understands design tradeoffs & 13 \\
\addlinespace
Process & 25 & Strong critical thinking in selecting design methods; prioritizes stakeholders & 13 \\
\addlinespace
Presentation & 20 & Extremely well-organized and well-coordinated; creates excitement & 10 \\
\addlinespace
Graphics & 10 & Original, compelling, effective; cites all external images & 5 \\

\end{longtable}
}

% ─── FDR Report Rubric ────────────────────────────────────────────
\subsubsection{Final Design Report Rubric}
\label{sssec:rubric-fdr-report}

{\small
\begin{longtable}{p{3.2cm} c p{3.5cm} p{3.5cm} p{2.8cm}}
\toprule
\textbf{Criterion} & \textbf{Pts} & \textbf{Exceeds} & \textbf{Meets / Adequate} & \textbf{Marginal} \\
\midrule
\endfirsthead
\toprule
\textbf{Criterion} & \textbf{Pts} & \textbf{Exceeds} & \textbf{Meets / Adequate} & \textbf{Marginal} \\
\midrule
\endhead
\bottomrule
\endfoot

Project Overview & 10 & Clear history and chronology; usable by future engineers & Generally clear but missing content & Limited documentation \\
\addlinespace
Design Approach & 10 & Problem-specific approach with correctly applied tools; meets constraints & Basic approach but lacking clarity & Vague or incomplete methodologies \\
\addlinespace
Engineering Analysis & 15 & Exceptionally detailed, error-free, no gaps; incorporates prior iteration feedback & Some gaps; questions unanswered & Not up to senior-level standard \\
\addlinespace
Final Solution & 30 & Fully described; testing demonstrates criteria met; supporting documentation complete & Adequate but incomplete documentation & Solution does not meet criteria \\
\addlinespace
Project Management & 5 & Accurate, up-to-date budget, WBS, schedule; clear path forward & Present but lacking detail & Substantial omissions \\
\addlinespace
Lessons Learned & 5 & Thoughtful discussion of deviations and key takeaways & Adequate but could be deeper & Conflicts with PA observations \\
\addlinespace
Technical Writing & 15 & Clear, concise, one voice & Several errors; some parts difficult to read & Many errors; not senior-level \\
\addlinespace
Spelling \& Grammar & 5 & Error-free & A few errors & Too many for senior writing \\
\addlinespace
Formatting \& Citations & 5 & IEEE citations; consistent formatting throughout & Some citation/formatting issues & No citations \\

\end{longtable}
}

% ─── FDR Presentation Rubric ─────────────────────────────────────
\subsubsection{Final Design Review Presentation Rubric}
\label{sssec:rubric-fdr-pres}

Identical structure to the CDR Presentation Rubric: Context (20\%), Content (25\%), Process (25\%), Presentation (20\%), Graphics (10\%). See Section~\ref{sssec:rubric-cdr-pres}.

% ─── Ethics Rubric ────────────────────────────────────────────────
\subsubsection{Ethics Assignment Rubric}
\label{sssec:rubric-ethics}

{\small
\begin{longtable}{p{3.5cm} c p{4cm} p{4cm}}
\toprule
\textbf{Criterion} & \textbf{Pts} & \textbf{Full Marks} & \textbf{Inadequate} \\
\midrule
\endfirsthead
\bottomrule
\endfoot

Intro to Ethical Consideration & 20 & Thoughtful introduction; clearly relates to project; explains why it is an ethical challenge & Limited relevance to project \\
\addlinespace
Affected Stakeholders & 15 & Thoughtful consideration of parties, interests, responsibilities & Limited discussion \\
\addlinespace
Alternative Courses of Action & 15 & Thoughtful alternatives with consequences to stakeholders & Limited discussion \\
\addlinespace
Ethics \& Liability Elaboration & 15 & Effectively uses Code of Ethics; discusses liability impacts & Minimal use of codes in context \\
\addlinespace
Rationale for Chosen Path & 20 & Clear, concise rationale based on ethical frameworks and values & Limited discussion \\
\addlinespace
Technical Writing & 10 & Clear, concise, no grammatical errors & Many errors; difficult to read \\
\addlinespace
Citations \& Formatting & 5 & IEEE format; all non-general-knowledge cited & No references to Codes of Ethics \\

\end{longtable}
}

% ─── Broader Impacts Rubric ──────────────────────────────────────
\subsubsection{Broader Impacts Essay Rubric}
\label{sssec:rubric-broader}

{\small
\begin{longtable}{p{3.5cm} c p{4cm} p{4cm}}
\toprule
\textbf{Criterion} & \textbf{Pts} & \textbf{Excellent} & \textbf{Needs Improvement / Poor} \\
\midrule
\endfirsthead
\bottomrule
\endfoot

Relevant & 25 & Addresses prompt in a controversial, insightful, or novel way & Does not understand intent of assignment \\
\addlinespace
Persuasive & 25 & Inspiring, thought-provoking; strong evidence of engineering education & Not overly creative; arguments weak \\
\addlinespace
Argumentative Essay Format & 10 & Well-organized; thesis, proponent/opponent arguments, strong conclusion & Reads more like informative essay or report \\
\addlinespace
Well Written & 20 & Compact, efficient, concise; could be considered for publication & Rambling, narrative, or noisy writing \\
\addlinespace
Spelling, Grammar, Punctuation & 10 & Practically no errors; strong evidence of proofreading & Pervasive errors; significant distraction \\
\addlinespace
References \& Citations & 10 & Thoroughly researched; IEEE format; scholarly sources & Little to no evidence of research \\

\end{longtable}
}

% ─── SoW / Project Plan Rubric ───────────────────────────────────
\subsubsection{Project Plan (Statement of Work) Rubric}
\label{sssec:rubric-sow}

{\small
\begin{longtable}{p{3.5cm} c p{4cm} p{4cm}}
\toprule
\textbf{Criterion} & \textbf{Pts} & \textbf{Exceeds} & \textbf{Marginal / No Evidence} \\
\midrule
\endfirsthead
\toprule
\textbf{Criterion} & \textbf{Pts} & \textbf{Exceeds} & \textbf{Marginal / No Evidence} \\
\midrule
\endhead
\bottomrule
\endfoot

Project Plan Summary & 2 & Professional, complete, concise & Missing \\
\addlinespace
Scope of Work & 5 & Clear, succinct problem statement; constraints well defined & Could be much improved \\
\addlinespace
Customer Needs & 5 & Clear, detailed list with importance ratings & Not provided \\
\addlinespace
Design Requirements & 10 & Tabulated (not bulleted); complete; linked to client needs & Incomplete and vague \\
\addlinespace
Stakeholders & 3 & Clear articulation with roles and timing & Insufficient detail \\
\addlinespace
Budget & 2 & Discusses allocation and fundraising plans & Not addressed \\
\addlinespace
Background Research & 10 & Demonstrates excellent understanding; existing solutions discussed & No evidence of understanding \\
\addlinespace
Eng. Standards & 9 & Identifies standards/regulations with detail on project applicability & One or no standard identified \\
\addlinespace
Material Handoff & 3 & Solid list of deliverables to client & Limited \\
\addlinespace
Performance Criteria & 3 & Thoughtful consideration of V\&V approach & Limited attempt \\
\addlinespace
Design Process & 5 & Thoughtful discussion of engineering design steps & Limited attempt \\
\addlinespace
Project Schedule & 5 & Detailed Gantt chart with discussion & Only reports/reviews; missing milestones \\
\addlinespace
Team Members & 3 & Table with names, skills, roles, contact info, advisors & Poorly constructed \\
\addlinespace
Management Tools & 3 & Strong management plan articulated & Not addressed \\
\addlinespace
Communication & 3 & Frequency, mediums, status report strategy & Missing information \\
\addlinespace
Technical Writing & 20 & Clear, concise, one voice, proofread & Lacks clarity; no attempt to edit \\
\addlinespace
Formatting \& Citations & 2 & IEEE citations; consistent formatting & No citations; poorly formatted \\

\end{longtable}
}

% ─── Design Showcase Poster Rubric ────────────────────────────────
\subsubsection{Design Showcase Poster Rubric}
\label{sssec:rubric-poster}

{\small
\begin{longtable}{p{4cm} c p{4.5cm} p{3.5cm}}
\toprule
\textbf{Criterion} & \textbf{Pts} & \textbf{Excellent} & \textbf{Poor} \\
\midrule
\endfirsthead
\bottomrule
\endfoot

Tables, Graphs, Figures & 5 & Correctly labeled, titled, explained; visuals convey information & Not labeled; visuals do not convey info \\
\addlinespace
Grammar \& Spelling & 5 & Free of all errors & Significant errors; poor editing \\
\addlinespace
Quality of Poster & 5 & Legible from 5+ ft; appropriate color/font; phrases not sentences & Not legible; poor color/font; text-heavy \\
\addlinespace
Design Process / Terminology & 5 & Clear presentation of process and final selection; appropriate terminology defined & No design process presented \\
\addlinespace
Logical Organization & 5 & Clear language and images in logical sequence; includes analysis, intro, conclusion, future work & Illogical sequence; little relevant info \\

\end{longtable}
}

% ─── Individual Performance Evaluation Rubric ─────────────────────
\subsubsection{Individual Performance Evaluation Rubric}
\label{sssec:rubric-ipe}

{\small
\begin{longtable}{p{3.5cm} c p{4cm} p{4cm}}
\toprule
\textbf{Criterion} & \textbf{Pts} & \textbf{Excellent} & \textbf{Poor / Unsatisfactory} \\
\midrule
\endfirsthead
\bottomrule
\endfoot

Technical Contribution & 40 & Produces excellent work; goes above and beyond; seeks help when appropriate; submits early for review & Poor quality; needs to be assigned work; leaves tasks for others; work often late \\
\addlinespace
Professional Contribution & 40 & Positive influence; completes all sprint tasks; helps others; demonstrates leadership; always on time; professional with all stakeholders & Unprofessional at times; does not complete sprint tasks; consistently late; unavailable outside class \\
\addlinespace
Assignment Submission & 20 & Fully addresses technical and professional contributions with clear supporting evidence & Minimally addresses contributions with no evidence \\

\end{longtable}
}

% ─── Next Steps Email Rubric ─────────────────────────────────────
\subsubsection{Next Steps Email Rubric (SDI)}
\label{sssec:rubric-next-steps}

{\small
\begin{longtable}{p{4cm} c p{5cm}}
\toprule
\textbf{Criterion} & \textbf{Pts} & \textbf{Full Marks} \\
\midrule
\endfirsthead
\bottomrule
\endfoot

PDR Action Items & 10 & Clear, complete list with status \\
\addlinespace
Project Budget \& Plan & 10 & Complete budget provided; original plan referenced; updates discussed \\
\addlinespace
Next Steps Discussion & 10 & Clear path forward; client expectations addressed; concerns noted \\
\addlinespace
Technical Writing & 10 & Clear, concise, one voice \\
\addlinespace
Spelling \& Grammar & 5 & No errors \\
\addlinespace
Formatting \& Citations & 5 & IEEE format; consistent formatting \\

\end{longtable}
}


% ══════════════════════════════════════════════════════════════════
\subsection{Grade Weights}
\label{subsec:grade-weights}
% ══════════════════════════════════════════════════════════════════

\subsubsection{EDNS 491 --- Senior Design I}

\begin{tabularx}{\textwidth}{Xr}
\toprule
\textbf{Assignment Group} & \textbf{Weight} \\
\midrule
Team Assignments --- Design Reports & 35\% \\
Individual Performance Evaluation & 20\% \\
Individual Writing Assignments & 10\% \\
Team Assessment --- Project Progress & 10\% \\
Team Assignments --- Design Reviews & 10\% \\
Prof.\ Dev.\ --- Engineering Standards & 5\% \\
Prof.\ Dev.\ --- New Employee Training (Complete All) & 3\% \\
Team --- Design Modules & 3\% \\
EHS Safety Training \& Access to Labriola Spaces & 2\% \\
Prof.\ Dev.\ --- Peer Evals, Surveys, Reflections & 2\% \\
\midrule
\textbf{Total} & \textbf{100\%} \\
\bottomrule
\end{tabularx}

\subsubsection{EDNS 492 --- Senior Design II}

\begin{tabularx}{\textwidth}{Xr}
\toprule
\textbf{Assignment Group} & \textbf{Weight} \\
\midrule
Performance Evaluations --- Individual & 25\% \\
Project Planning \& Completion --- Team & 25\% \\
Individual Writing Assignment & 10\% \\
Design Reviews --- Team & 10\% \\
Critical Design Report --- Team & 10\% \\
Capstone Design Showcase --- Team & 10\% \\
Professional Development --- Individual & 5\% \\
Final Design Report --- Team & 5\% \\
Final Checkout & 0\% \\
\midrule
\textbf{Total} & \textbf{100\%} \\
\bottomrule
\end{tabularx}

\begin{protip}
In SDI, the biggest grade drivers are Design Reports (35\%) and Individual Performance Evaluation (20\%). In SDII, Performance Evaluations (25\%) and Project Planning \& Completion (25\%) dominate. Prioritize accordingly.
\end{protip}
